\section{Dispatch buffer as glue between Postgres scan and alignment pipeline}
\label{sec:match-placeholder-injection}

I still need to explain how exactly \texttt{CustomBioseqScan} works with worker-architecture shown in Figure~\ref{fig:sequence-scan-lifecycle}.
The main abstraction here is the dispatch buffer which is reposible for transferring data between Postgres and workers.

Figure~\ref{fig:execcustomscan-dispatch} shows how Panel~(a) of Figure~\ref{fig:sequence-scan-lifecycle} looks in the context of Postgres scan.
\texttt{BioseqScan} asks for the next row from TAM. This row is placed into dispatch buffer (locator only) and dispatch buffer holds all the structures presented before~\ref{fig:sequence-scan-lifecycle} - scanners.
The dispatch buffer code then handles all view copies to the workers.

\begin{figure}[htbp]
\centering
\begin{tikzpicture}[x=1cm,y=1cm,font=\scriptsize,
  pg/.style={scan box,minimum width=3.55cm,minimum height=.56cm,
    font=\scriptsize,align=center,inner sep=2pt,outer sep=0.5pt},
  hook/.style={scan hit,minimum width=4.05cm,minimum height=.56cm,
    font=\scriptsize,align=center,inner sep=2pt,outer sep=0.5pt},
  tam/.style={scan output,minimum width=3.85cm,minimum height=.56cm,
    font=\scriptsize,align=center,inner sep=2pt,outer sep=0.5pt},
  stagelab/.style={anchor=north west,font=\scriptsize\bfseries,
    text=black!45,inner sep=0pt},
  note/.style={font=\scriptsize,text=black!60,inner sep=1pt,align=center},
  callarr/.style={scan arrow},
  altarr/.style={scan arrow,draw=scanBlue}]
% three code regions
\fill[scanOrange!8,rounded corners=1.2pt]
  (0.06,-1.70) rectangle (4.78,9.88);
\node[stagelab] at (0.16,9.80) {BioseqScan};
\fill[scanBlue!9,rounded corners=1.2pt]
  (4.90,-1.70) rectangle (9.78,9.88);
\node[stagelab] at (5.00,9.80) {Dispatch buffer code};
\fill[scanGreen!8,rounded corners=1.2pt]
  (9.90,-1.70) rectangle (15.50,9.88);
\node[stagelab] at (10.00,9.80) {TAM};
% BioseqScan
\node[pg] (exec) at (2.42,9.20) {\texttt{ExecCustomScan}};
\node[hook] (locator) at (2.42,8.00) {locator row};
\node[pg] (retnow) at (2.42,2.30) {return one row};
% Dispatch buffer
\node[draw=scanOrange,fill=scanOrange!12,rounded corners=2.5pt,
  minimum width=4.37cm,minimum height=1.50cm,anchor=center] (bufbox) at (7.34,6.30) {};
\node at (7.34,6.82) {locator slots};
\foreach \i in {0,1,2,3} {
  \draw[draw=scanOrange,fill=white]
    (5.55+0.90*\i,5.78) rectangle ++(.78,.48);
  \node at (5.94+0.90*\i,6.02) {\tiny loc};
}
\node[draw=black!40,rounded corners=2pt,dashed,
  minimum width=4.37cm,minimum height=2.10cm,anchor=center] (scanframe) at (7.34,4.12) {};
\node[anchor=north west,font=\scriptsize\itshape,text=black!70]
  at (5.25,5.10) {session + scanners};
\foreach \i in {0,1} {
  \draw[draw=black!45,fill=white]
    (5.30,4.45-0.72*\i) rectangle (9.38,4.82-0.72*\i);
  \foreach \j in {0,...,9}
    \draw[draw=black!40,fill=white]
      (5.42+0.38*\j,4.52-0.72*\i) rectangle ++(.32,.24);
}
\foreach \j in {0,1,4}
  \draw[draw=scanBlue,fill=scanBlue!35]
    (5.42+0.38*\j,4.52) rectangle ++(.32,.24);
\node[draw=scanGreen,fill=scanGreen!14,
  minimum width=1.18cm,minimum height=.72cm] (w1) at (5.88,0.55) {worker 1};
\node[draw=scanGreen,fill=scanGreen!14,
  minimum width=1.18cm,minimum height=.72cm] (w2) at (7.34,0.55) {worker 2};
\node[draw=scanGreen,fill=scanGreen!14,
  minimum width=1.18cm,minimum height=.72cm] (w3) at (8.80,0.55) {worker 3};
% TAM
\node[tam] (usertam) at (12.24,9.20) {user table};
\node[note] at (13.48,8.58) {locator + metadata};
\node[tam] (colcache) at (12.24,4.12) {column cache};
\node[note] at (11.18,2.82) {load stripe};
\node[draw=scanGreen,fill=scanGreen!9,rounded corners=2pt,
  minimum width=4.08cm,minimum height=1.70cm,anchor=center] (fragbox) at (12.24,1.00) {};
\node[font=\scriptsize] at (12.24,1.62) {fragment sub-table};
\foreach \i in {0,1,2,3} {
  \draw[draw=scanGreen,fill=scanGreen!22]
    (10.45+0.90*\i,0.32) rectangle (11.20+0.90*\i,1.28);
  \node[font=\tiny] at (10.825+0.90*\i,0.80) {f\i};
}
% ExecCustomScan asks the user-table TAM for the next locator
\draw[callarr] (exec)--(usertam);
% locator comes back: drop below the user table, then a long run into the east side
\draw[callarr,rounded corners=2.5pt]
  (usertam.south)--++(0,-0.92)--(locator.east);
% place into buffer from the middle of the locator's south side
\draw[callarr,rounded corners=2.5pt]
  (locator.south)--++(0,-1.42)--(bufbox.west);
\node[note,align=center] at (4.05,6.55) {place into\\buffer};
% not-full loop: long last segment into the east side, not a tight top corner
\draw[callarr,dashed,rounded corners=2.5pt]
  (bufbox.east)--(14.85,6.30)--(14.85,9.06)--([yshift=-4pt]usertam.east);
\node[note,align=center] at (13.12,6.85)
  {fetch next row\\if batch not full};
% full buffer opens scanners on the backend thread
\draw[callarr] (bufbox.south)--(scanframe.north);
\node[note] at (6.18,5.34) {run next batch};
% scanners load fragments. workers never touch TAM or the cache
\draw[callarr] (scanframe.east)--(colcache.west);
\node[note] at (9.82,4.42) {miss};
\draw[callarr] (colcache.south)--(fragbox.north);
\draw[callarr,dashed]
  ([yshift=-8pt]colcache.west)--([yshift=-8pt]scanframe.east);
\node[note] at (9.82,3.78) {copy};
\coordinate (viewbus) at (7.34,1.90);
\draw[rounded corners=2pt,thick] (scanframe.south)--(viewbus);
\draw[callarr,rounded corners=2.5pt] (viewbus)-|(w1.north);
\draw[callarr,rounded corners=2.5pt] (viewbus)--(w2.north);
\draw[callarr,rounded corners=2.5pt] (viewbus)-|(w3.north);
\node[note] at (8.72,2.22) {provide view copy};
% drain hits enter the east side of the return box, not along its bottom
\draw[callarr,rounded corners=2.5pt]
  (w1.west)--++(-0.45,0)|-(retnow.east);
\node[note] at (3.55,1.15) {drain hits};
\draw[callarr,dashed,rounded corners=2.5pt]
  (retnow.west)--++(-0.55,0)|-(exec.west);
\node[note,align=left,anchor=west] at (0.52,5.85) {Postgres fetches\\next tuple};
% no search: leave the locator to the right, hydrate via cache, return into south
\draw[altarr,rounded corners=2.5pt]
  (locator.south east)--++(0.40,0)--++(0,-0.28)--(15.15,7.44)
  --(15.15,4.12)--(colcache.east);
\draw[altarr,rounded corners=2.5pt]
  ([xshift=20pt]colcache.south)--++(0,-0.48)--(15.15,3.36)
  --(15.15,-1.12)--(2.42,-1.12)--(retnow.south);
\node[note,text=scanBlue,align=center] at (7.34,-1.40)
  {no search: hydrate via cache, then return};
\end{tikzpicture}
\caption{Inside \texttt{ExecCustomScan}, the following things happen:
\texttt{BioseqScan} holds dispatch buffer and fetches a locator handles it to the buffer logic and retrieves results from it.
The dispatch buffer stores locators, then builds a session that holds scanners.
A scanner miss reads the fragment sub-table through the TAM column cache and copies the fragment into a slot.
Workers receive a view copy from the scanners.
Black arrows show the normal alignmen execution path, however it might happen that we don't perform alignment and just want to return the sequences itself.
In that case the flow uses blue arrows (there's no \texttt{mmsearch} in the query) and \texttt{BioseqScan} hydrates the sequence row through cache and just returns it without interacting with dispatch buffer.
Dashed arrows repesent loops here.}
\label{fig:execcustomscan-dispatch}
\end{figure}

The collect loop is the gap between batch~1 and batch~2 in panel~(d) of Figure~\ref{fig:sequence-scan-lifecycle}.
\texttt{BioseqScan} may call the user-table TAM many times before it returns to fill in the dispatch buffer.
The backend thread that owns the scanners performs fragment reads, while workers only see the view copy described in panel~(b).



Assigning the search to a custom scan solves the missing execution order.
The owner selects a hit before PostgreSQL evaluates the expressions that
produce its output row. Those expressions read fields already computed by
the aligner. They do not call it again.

The argument in \texttt{bioseq.dp\_score(t.seq)} is used to identify the
search target. The score is not stored in \texttt{t.seq}. This distinction
also explains why the same expression must not silently read a hit belonging
to another search. For a supported transformed target, such as a slice, the
corresponding helper can identify that same transformed expression.

\begin{figure}[htbp]
\centering
\begin{tikzpicture}[font=\footnotesize]
  \node[anchor=west,font=\bfseries] at (0,4.5) {Two hits from one target row};
  \node[scan input,minimum width=2.4cm] (row) at (1.25,3.3) {target id 27};
  \node[scan box,minimum width=3.1cm,minimum height=.9cm] (hit0) at (5.0,3.8)
    {hit 0\\score 42, start 100};
  \node[scan hit,minimum width=3.1cm,minimum height=.9cm,very thick] (hit1) at (5.0,2.4)
    {hit 1\\score 38, start 800};
  \draw[scan arrow] (row.east) -- (hit0.west);
  \draw[scan arrow,scanOrange] (row.east) -- (hit1.west);
  \node[anchor=west,text=scanOrange] at (3.3,1.65) {current hit};
  \node[scan output,minimum width=2.0cm] (score) at (9.8,3.8)
    {score: 38};
  \node[scan output,minimum width=2.0cm] (start) at (9.8,2.4)
    {start: 800};
  \node[scan output,minimum width=2.0cm] (cigar) at (9.8,1.0)
    {CIGAR of hit 1};
  \draw[scan arrow,scanOrange] (hit1.east) to[out=20,in=180] (score.west);
  \draw[scan arrow,scanOrange] (hit1.east) -- (start.west);
  \draw[scan arrow,scanOrange] (hit1.east) to[out=-20,in=180] (cigar.west);
  \node[anchor=west,text=black!65] at (0,.2)
    {Before the next output row, the owner selects the next hit and updates all fields.};
\end{tikzpicture}
\caption{The current hit is the common source of score, position, and CIGAR.
The numbers are illustrative. Changing the current hit changes all projected
fields together while the target identifier stays the same.}
\label{fig:scan-current-hit}
\end{figure}

\subsection{Typed values in the expression tree}

PostgreSQL evaluates output expressions through an expression tree. Before
returning its first hit, the executor walks this tree and replaces the helpers
belonging to the owner with typed \texttt{Const} nodes. It retains their
addresses. Before projecting a hit, it writes the hit's values and null flags
into these nodes and rebuilds the output-expression evaluator so that it
reads the changed fields.

\begin{itemize}
  \item \textbf{Identity answers which search.} The argument's target
    binding links a helper to the right owner. A function name alone cannot
    distinguish two searches in one statement.
  \item \textbf{Type answers how to read the field.} Coordinates use
    \texttt{BIGINT}, the DP score uses \texttt{INT4}, Boolean properties use
    \texttt{BOOLEAN}, and CIGAR uses \texttt{TEXT}.
  \item \textbf{The current hit answers which value.} Every field is filled
    from the same hit record before the output row is returned.
\end{itemize}

Although PostgreSQL calls the node a \texttt{Const}, this use is local to the
executor. The value changes between output rows. The replacement happens
after planning, so the planner does not fold a current hit's score into a
reusable plan. No session-wide ``last hit'' variable is needed.
ParadeDB's base scan uses the same kind of changing score node and rebuilds
the projection after updating it.
\coderef{paradedb/pg_search/src/postgres/customscan/basescan/mod.rs}
  {ParadeDB current-result score}{exec_custom_scan}{1784}

\begin{algorithm}[htbp]
\caption{Project one hit without repeating alignment}
\label{alg:scan-project-hit}
\small
\begin{enumerate}
  \item Restore the ordinary input columns for this hit.
  \item For every requested helper, write its field from this hit into the
    corresponding typed value node.
  \item Rebuild the evaluator, then evaluate the output expressions.
  \item Return the resulting SQL row. Keep the next hit for the next request.
\end{enumerate}
\end{algorithm}

The ordinary helper bodies are not a uniform substitute for this mechanism.
Some helpers report an invalid use, while some generic numeric helpers return
a special fallback value when called outside a supported owner. Such values
are not computed alignment results.

\subsection{Placement during planning is a separate problem}

Joins can move expressions between plan levels. An outer join can also add a
row of nulls for a missing input. PostgreSQL's \texttt{PlaceHolderVar} records
where an expression must be evaluated and which outer joins can make its
result null.
\coderef{postgres17/src/include/nodes/pathnodes.h}
  {PostgreSQL expression placement}{PlaceHolderVar}{2750}
BioSeqDB uses this planning mechanism where needed to keep a helper attached
to the right input.
ParadeDB uses this mechanism to keep score expressions at their search
input while planning joins and aggregates.
\coderef{paradedb/pg_search/src/postgres/customscan/projections.rs}
  {ParadeDB score placement}{placeholder_support}{314}

\texttt{PlaceHolderVar} and the changing \texttt{Const} nodes do different
jobs. The first preserves expression placement while the planner builds and
rewrites the tree. The second supplies a field of the current hit during
execution. Neither permits unsupported join shapes to bypass the owner's
validation.

\subsection{Requested fields can require more alignment work}

The owner also records which fields the SQL statement needs. For example,
CIGAR output requires a base-level edit path and a formatted string. A
coordinate-only query does not need that string. The requested fields are
combined with the selected alignment options when the executor starts.

There is a practical limit in the current implementation: requesting CIGAR
can enable further alignment work, and that path drops hits for which no
usable CIGAR is produced. Therefore, adding CIGAR to a query is not guaranteed
to leave its hit count unchanged. This behavior must be considered when
comparing counts from queries with different output fields.
